\begin{claim}{D.66}{Proved}{The exact norm in a scalar cyclic torus cover}
Let $a$ be an actual element of order $d=3$ or $4$ in the
full once-punctured-torus loop-outer group, and put $p=3$
or $2$, respectively. At every scalar coordinate retained
by D.54--D.55, write
\[
 A_q(a)=uI,\qquad
 \begin{cases}
  u=\lambda_q^2,\quad \lambda_q\ne1,&d=3,\\
  u^2=-1,\quad \lambda_q=-1,&d=4.
 \end{cases}                                                   \tag{D.66a}
\]
Thus $q\equiv1\pmod d$, and the reduction $\bar u\in\F_q^\times$
has exact order $d$. For every nonzero character
$\chi:F=\widehat F(x,y)\to C_q$, its kernel $N$ is the full
completion of the actual degree-$q$ cover of type $(1,q)$.
The complete primary homology comparison is
\[
 H_1(N,\Z_p)
   \cong (\Z_p^2,A_p(a))\ \oplus\
          \Z_p[C_d]^{(q-1)/d}.                              \tag{D.66b}
\]
Here the action is the one induced by a labelled cusp-normalizer
element above $a$; the two displayed summands are the actual
deck-average handle summand and boundary summand.

Label the boundary classes by $b_i$, $i\in\F_q$, with the
distinguished original cusp labelled $0$. The boundary module
has the exact relation $\sum_i b_i=0$, and the action is
$b_i\mapsto b_{\bar u i}$. If $R$ contains one representative
of each nonzero $\langle\bar u\rangle$-orbit, then
\[
 w=-\sum_{r\in R}b_r,\qquad
 (1+T+\cdots+T^{d-1})w=b_0,                              \tag{D.66c}
\]
where $T$ denotes this action. In particular the distinguished
unit cusp-power class is an integral norm, at every primary
precision. More fully, for $M=H_1(N,\Z_p)$,
\[
 H^2(C_d,M)=0,\qquad H^1(C_d,M)\cong C_p.                \tag{D.66d}
\]
The actual cyclic extension over this cover, after quotienting
its kernel to $M$, has the explicit splitting given below.
These statements identify exactly why the invariant-character
norm contradiction of D.54--D.55 does not extend to the scalar
coordinates (D.66a). They neither construct an action having
such a coordinate nor decide the remaining scalar multiplier
or conjugacy questions.
\end{claim}
\begin{proof}
\emph{The actual extension and its finite quotient.}
Choose an actual representative $\alpha$ of order $d$,
using D.27 for $d=3$ and D.46 for $d=4$, and set
\[
 E=F\rtimes_\alpha\langle s\rangle,\quad c=[x,y],\quad
 L=\langle c\rangle .
\]
The chosen $s$ has quotient label $a$. Choose a cusp-normalizer
element $v$ with the same quotient label. The exact relations are
\[
 vcv^{-1}=c^\lambda,\qquad v^d=c^\kappa,\qquad
             \lambda_p=1,\quad \kappa_p\in\Z_p^\times.       \tag{D.66e}
\]
For $d=3$ these are D.48. For $d=4$, the element $v^2$
lies over the positive involution $a^2$. Its square exponent
is a two-adic unit by D.55's full cusp-normalizer calculation.
Also $\lambda_2=1$: the order-four matrix on $\Z_2^2$ has
square $-I$, characteristic polynomial $X^2+1$, and determinant
one. The peripheral multiplier equals that determinant.

Let $\bar u$ be as in (D.66a), and write $C_q$ additively.
Both conjugation by $s$ and by $v$ satisfy
\[
             \chi(vhv^{-1})=\bar u\,\chi(h)
                         \qquad(h\in F).                    \tag{D.66f}
\]
Indeed they induce the same homology map, and every character
factors through the actual abelianization of $F$.
Thus $N=\ker\chi$ is normal in $E$. Since $c\in N$,
the class of $v$ in $E/N$ has exact order $d$.
Choose a discrete word $f$ with $\chi(f)=1$. The two mutually
inverse finite quotient maps are
\[
 \begin{aligned}
 C_q\rtimes_{\bar u} C_d&\longrightarrow E/N,
             &(i,j)&\longmapsto f^i v^jN,\\
 E/N&\longrightarrow C_q\rtimes_{\bar u} C_d,
             &hv^jN&\longmapsto(\chi(h),j),
                    \quad 0\leq j<d.
 \end{aligned}                                               \tag{D.66g}
\]
The multiplication on the left is
$(i,j)(k,l)=(i+\bar u^j k,j+l)$; the relation $v^d\in N$
and (D.66f) verify both maps and their inverse identities.

Put $E_N=N\langle v\rangle$, the inverse image of
$\{0\}\rtimes C_d$ in this exact quotient. It has index
$q$ in $E$, with kernel $N$ and quotient $C_d$.
For completeness it has an actual order-$d$ section.
Write the image of the original $s$ as $(j,1)$ and choose
the integer residue
\[
                    k=\frac{j}{\bar u-1}\quad\text{in }\F_q.
\]
The denominator is nonzero. Conjugation by $(k,0)$ sends
$(j,1)$ to $(j+(1-\bar u)k,1)=(0,1)$. Consequently
$s_0=f^k s f^{-k}$ belongs to $E_N$, has exact order $d$,
and gives $E_N=N\rtimes\langle s_0\rangle$ with the
specified quotient label. No character extension
$E\to C_q$ with trivial action is being substituted for
(D.66g).

\emph{The actual cover and the boundary labels.}
The restriction of $\chi$ to the dense discrete free group
is onto. Its punctured-torus cover has degree $q$ and
trivial cusp monodromy, hence $q$ degree-one boundary
lifts. Euler characteristic $-q$ gives genus one.
The finite-index completion map identifies its full
fundamental-group completion with precisely $N$.
The chosen boundary representatives are
\[
           c_i=f^i c f^{-i}\in N\qquad(0\leq i<q).
                                                               \tag{D.66h}
\]
They have distinct cover labels. In homology their classes
$b_i$ have exactly the single relation $\sum_i b_i=0$.
For $j$ the representative of $\bar u i$ modulo $q$, put
$n_i=(vf^iv^{-1})f^{-j}\in N$, by (D.66f). Then the
actual conjugacy equation is
\[
                 vc_iv^{-1}=n_i c_j^\lambda n_i^{-1}.     \tag{D.66i}
\]
In $p$-primary homology $\lambda_p=1$, so its action is
exactly $T b_i=b_{\bar u i}$. The label zero is fixed,
and every nonzero label has orbit length $d$.

\emph{Deck averaging and both inverse boundary-module maps.}
Let $D=F/N=C_q$ act on $M=H_1(N,\Z_p)$ by conjugation
of deck lifts, and set
\[
                  e=\frac1q\sum_{h\in D}h.                \tag{D.66j}
\]
The prime $q$ differs from $p$. Although $T$ need not
commute with individual deck transformations, it permutes
them by $h\mapsto h^{\bar u}$; hence it commutes with $e$.
Filling all $q$ boundaries gives the actual compact torus
cover associated to the lattice
$\Lambda=\ker(\Z^2\to C_q)$, with covering map
$\R^2/\Lambda\to\R^2/\Z^2$. Its deck maps are translations,
so they act trivially on the filled homology. The boundary
module is
\[
 B_\partial=\Z_p\{b_i:i\in\F_q\}/\Z_p\Bigl(\sum_i b_i\Bigr).
                                                               \tag{D.66k}
\]
It has zero deck average, since $q$ is a unit. The boundary
sequence therefore gives the exact identifications
\[
 (1-e)M=B_\partial,\qquad
 eM\xrightarrow{\ \cong\ }H_1(\R^2/\Lambda,\Z_p)
       \xrightarrow{\ \cong\ }H_1(\R^2/\Z^2,\Z_p).       \tag{D.66l}
\]
The second map is induced by the index-$q$ lattice inclusion;
its determinant is a $p$-adic unit. Both maps intertwine
the induced homology actions, so $T|_{eM}=A_p(a)$ under
these specified identifications.

Write $m=(q-1)/d$. With $R$ as in the statement, define
\[
 \begin{aligned}
 \Phi:\bigoplus_{r\in R}\Z_p[C_d]e_r&\longrightarrow B_\partial,
                    &T^j e_r&\longmapsto b_{\bar u^j r},\\
 \Psi:B_\partial&\longrightarrow
                         \bigoplus_{r\in R}\Z_p[C_d]e_r,
                    &b_{\bar u^j r}&\longmapsto T^j e_r,\\
                    &&b_0&\longmapsto
                         -\sum_{r\in R}\sum_{j=0}^{d-1}T^j e_r.
 \end{aligned}                                               \tag{D.66m}
\]
The displayed image of the total boundary relation is zero.
It also makes $\Psi$ equivariant at the fixed label zero.
Every nonzero label has a unique expression $\bar u^j r$,
so these maps are mutually inverse on all their generators.
This proves the free cyclic-module assertion in (D.66b).
Applying $\Phi$ to the negative sum of the $e_r$ proves
(D.66c) integrally.

\emph{The exact cyclic cohomology and abelianized extension.}
On the handle summand $H=eM$, D.48 and D.55 give respectively
\[
 \begin{array}{c|c|c|c}
 d&p&\text{identity on }H&\det(T-I)|_H\\ \hline
 3&3&T^2+T+I=0&3\\
 4&2&T^2=-I&2
 \end{array}                                                   \tag{D.66n}
\]
Thus the cyclic norm $\mathcal N=1+T+\cdots+T^{d-1}$ is
zero on $H$, while $H^T=0$. The cokernel of $T-I$ on
$H\cong\Z_p^2$ is $C_p$: its determinant has valuation one,
so its two invariant factors have valuations zero and one.
On each free summand $\Z_p[C_d]$, the invariants are
exactly the multiples of the sum of its $d$ basis elements,
which are exactly its norm image. The kernel of the norm
is the augmentation ideal, the image of $T-I$.
The periodic cyclic cochain complex consequently gives
\[
 M^T/\mathcal N M=0,\qquad
 \ker(\mathcal N)/(T-I)M\cong C_p,
\]
which is precisely (D.66d).

Here is an explicit splitting retaining the cusp exponent.
Let $K=\ker(N\to M)$ and $P=E_N/K$. This is an actual
quotient, since $K$ is characteristic in $N$. Write the
kernel $M$ additively, and still denote the image of $v$
by $v$. Its relations are
\[
           vmv^{-1}=Tm,\qquad v^d=\kappa_p b_0.
                                                               \tag{D.66o}
\]
They are complete: every element of $P$ has a unique
normal form $m v^j$, $0\leq j<d$, and these relations
give its multiplication. By (D.66c), the element
\[
                    t=(-\kappa_p w)v
\]
has $t^d=0$ in the additive kernel, and has exact order
$d$ because of its quotient label. The inverse
isomorphisms are
\[
 \begin{aligned}
 M\rtimes_T C_d&\longrightarrow P,&(m,j)&\longmapsto m t^j,\\
 P&\longrightarrow M\rtimes_T C_d,&m&\longmapsto(m,0),\\
 &&v&\longmapsto(\kappa_p w,1).
 \end{aligned}                                               \tag{D.66p}
\]
The norm identity verifies the image of $v^d$, and
the two composites fix the indicated kernel and generator.
Reducing these maps modulo every $p^n$ gives the same
splitting at that exact precision. All module solutions
of $\mathcal N z=-\kappa_p b_0$ form
$-\kappa_p w+\ker\mathcal N$.

The actual section $s_0$ constructed above has one such
vector relative to $v$, but this calculation does not
assert that every vector in that coset lifts to a
finite-order element before the quotient by $K$.
In the invariant-character proof all boundary labels
were fixed and the primary norm vanished modulo $p$.
Here the explicit free boundary summands (D.66m) give
the precise integral preimage for the same distinguished
cusp class. Thus there is no remaining norm obstruction
from this cover module, without any assertion that a
scalar-coordinate full action exists.

The executable certificate
\texttt{code/part\_c/torus\_scalar\_cover\_norm\_d66.py}
checks the finite quotient conjugation, the two inverse
boundary-module maps, the deck conjugacy and averaging,
the integral norm preimage, and the handle matrices
for the first admissible cubic and quartic primes.
The arbitrary-prime assertions follow from the explicit
maps above, not from this finite enumeration.
\end{proof}
